\section*{How to read the construction}
\addcontentsline{toc}{section}{How to read the construction}

The starting question is elementary: why do two finite bases have the same
number of vectors? One answer exchanges vectors by triangular substitutions.
We begin by explaining those substitutions as invertible maps, and the
familiar algebraic maneuvers as commutative diagrams. We then ask for a
construction that can compare layers even when there are no vectors and
no addition of maps.

The answer uses five ingredients with distinct jobs. Finite limits impose
compatible equations. Regular images forget intermediate witnesses in a
controlled way. Exactness turns an equivalence relation into a quotient.
The Mal'cev condition makes every compatible relation difunctional:
three related corners force the fourth. Pointedness selects the zero class
needed to organize a filtration. These hypotheses will be defined and used
one at a time; ``exact'' here means Barr-exact, not Quillen-exact.

The recurring construction is
\[
 \text{one correspondence}
 \quad\longmapsto\quad
 \text{two ambiguity relations}
 \quad\longmapsto\quad
 \text{one canonically identified quotient}.
\]
Applying it at every crossing of two filtrations produces the refinement
rectangle. Jordan--H\"older follows when simple factors prevent a
nontrivial further decomposition. Groups and modules are verified
instances of this one proof. We do not assume a prior dimension theorem,
an invariant length, or a Grothendieck group to justify it.

For a first reading focused on that result, read
Sections~\ref{f:begin}--\ref{e:section}. The relation notation requires
some attention, but the convention never changes: $SR$ means first $R$,
then $S$. Each abstract formula is accompanied by its meaning and its
typed maps. The elementary algebra supplies the concrete models that make
the common construction intelligible.

For the spectral viewpoint, continue through the module toolbox and the
comparison pages. The groups-with-operators calculation works with actual
nonabelian quotients because its differential images have proved
normality. The subsequent module chapter establishes the more general
two-term filtered-map calculus. Neither calculation is used to patch a
gap in the categorical refinement proof.

For reconstruction, follow the next four chapters. Simple factors give
the universal additive record, but cannot distinguish all extensions.
Explicit attachment parameters, maps from small probes, and counts of
flags recover progressively different information. The nilpotent
reconstruction theorem is proved by constructing chain bases; it is not
inferred from a classification theorem left unstated.

The final part changes the kind of quotient being taken. A homotopy
quotient retains automorphisms and coherent transport that an orbit set
forgets. Its descent theorem, coefficient categories, and derived
extension and cotangent constructions offer additional capabilities.
There the foundational equivalences from higher category theory are
\emph{explicitly quoted inputs}. Our deductions from them are proved,
with their assumptions and maps. This makes the logical boundary visible:
the text is self-contained for its ordinary categorical comparison,
relative to the categorical axioms it defines; it is not a reconstruction
of all of higher topos theory from set theory.

No claim of historical novelty is made. The purpose is a connected
exposition in which constructions retain, forget, and recover specified
data. The dependency record at the end separates theorems proved here,
framework results imported from the literature, illustrative calculations,
and the limits of each conclusion.
